\section{Conclusion and Next Steps}
\label{sec:conclusion}

\subsection{Summary of Contributions}

This technical report presents IsalSR, a complete implementation of
isomorphism-invariant expression representation for symbolic regression.
The key contributions are:

\begin{enumerate}
  \item \textbf{Two-tier instruction set} ($\Sigma_{\mathrm{SR}}$) that
    encodes labeled DAGs as instruction strings with 31 tokens covering
    12 operation types (configurable per experiment).

  \item \textbf{Canonical string computation} ($w^*_D$) that is a complete
    labeled-DAG invariant, collapsing $O(k!)$ equivalent representations
    into one. Computed from $x_1$ only (no starting-node iteration) with
    6-component structural tuple pruning.

  \item \textbf{Complete software library} with 286 verified tests, strict
    type checking, zero-dependency core, and proposal-alignment validation
    across 7 development phases.

  \item \textbf{Experimental infrastructure} for validating the $O(k!)$
    reduction claim on standard benchmarks (Nguyen, Feynman, SRBench)
    with a paired WITH vs.\ WITHOUT canonicalization design.
\end{enumerate}

\subsection{Limitations}

\begin{itemize}[nosep]
  \item \textbf{Canonical computation cost}: The exhaustive backtracking is
    exponential in the number of V/v branch points. The pruned variant
    mitigates this but may still be prohibitive for DAGs with $k > 15$.
  \item \textbf{Binary operation ordering}: Non-commutative operations
    (SUB, DIV, POW) use sorted-by-node-ID input ordering, which must be
    consistent between canonicalization and evaluation.
  \item \textbf{Full-scale experiments pending}: The preliminary results
    are from small local runs. Publication-ready experiments require
    larger $N$, more runs, and statistical analysis.
\end{itemize}

\subsection{Roadmap to IEEE TPAMI Submission}

\begin{enumerate}[nosep]
  \item Run full search space analysis (10,000+ strings, all 12+10 benchmarks).
  \item Run paired WITH/WITHOUT experiments on random search, hill climbing,
    and evolutionary population.
  \item Compute statistics: mean $R^2$, max $R^2$, IQR, reduction factor.
  \item Generate publication figures using \texttt{plotting\_styles.py}.
  \item Write paper following IEEE TPAMI format.
  \item Submit arXiv preprint, then journal submission.
\end{enumerate}

\subsection{Target Venue}

\textbf{IEEE Transactions on Pattern Analysis and Machine Intelligence (TPAMI)}.
The contribution is a \textit{representation} result with mathematical proof
and empirical validation, not a new SR algorithm. This positions the work
as a foundational contribution to the field.
